\section{Ground spaces of heterogeneous open strips}%
\label{sec:peps_g_injective_strip_parent_ground_space}

\begin{definition}[Adjacent-pair conditions on an open strip]%
    \label{def:peps_g_injective_strip_parent_space}
    \lean{TNLean.PEPS.StripBlock,
        TNLean.PEPS.stripConfig,
        TNLean.PEPS.finite_stripConfig,
        TNLean.PEPS.stripTensor,
        TNLean.PEPS.stripFirstPairSlice,
        TNLean.PEPS.stripTailSlice,
        TNLean.PEPS.stripParentGroundSpace,
        TNLean.PEPS.mem_stripParentGroundSpace_cons_cons_cons_iff}
    \leanok
    A strip consists of an ordered list of tensor blocks
    $A_1,\ldots,A_n$ on a common finite-dimensional virtual space
    $\mathcal V$. Each block has its own finite physical alphabet
    $I_r$, and a physical configuration is $(i_1,\ldots,i_n)$.
    Its full boundary map is
    \[
        \mathcal P(A_1\cdots A_n)X
        =\bigl(\operatorname{tr}
        [A_1^{i_1}\cdots A_n^{i_n}X]\bigr)_{i_1,\ldots,i_n}.
    \]
    The common adjacent-pair space $\mathcal K_n$ consists of all
    physical vectors such that, after fixing every physical index
    outside a pair $r,r+1$, the remaining two-block vector belongs
    to $\operatorname{ran}\mathcal P(A_rA_{r+1})$.
    Equivalently, impose the first-pair condition on every remaining
    configuration and then impose the same adjacent conditions on
    every first-index-fixed tail. For $n=2$ there is just one pair
    condition; for $n\le1$ there are no conditions.
    These are the actual local parent conditions used in the
    intersection iteration of \cite[Theorem~4.11]{Schuch2010PEPS}.
\end{definition}

\begin{theorem}[Heterogeneous open-strip ground space]%
    \label{thm:peps_g_injective_strip_parent_space}
    \lean{TNLean.PEPS.isGInjective_stripTensor,
        TNLean.PEPS.stripParentGroundSpace_eq_range}
    \leanok
    \uses{def:peps_g_injective_strip_parent_space,def:peps_mps_g_injective}
    Let a finite group act on $\mathcal V$ through a common
    representation $\rho$. Suppose every block is $G$-injective
    and $n\ge2$. Then
    \[
        \mathcal K_n=\operatorname{ran}\mathcal P(A_1\cdots A_n).
    \]
    No equality of local intersections is a hypothesis. The blocks
    and their physical alphabets may vary along the strip.
    This is the open-boundary intersection iteration in
    \cite[Theorems~4.11 and~5.4]{Schuch2010PEPS}. A PEPS graph
    application additionally requires identifying its regional
    contractions with these operator-valued strip tensors.
\end{theorem}

\begin{proof}\leanok
    \uses{thm:peps_g_injective_strip_intersection,
        thm:peps_mps_g_injective_concatenation}
    For two blocks, the statement is the single pair condition.
    For a longer strip, apply induction to every first-index-fixed
    tail. The common conditions then become the first-pair space
    intersected with the full tail boundary space. Block the
    nonempty suffix after the second site into a tensor $C$.
    Concatenation makes $C$ $G$-injective. The three-block
    intersection theorem for $A_1,A_2,C$ gives precisely the
    full-strip boundary range.
\end{proof}

\begin{corollary}[Unique invariant boundary and ground-space dimension]%
    \label{cor:peps_g_injective_strip_parent_boundary}
    \lean{TNLean.PEPS.existsUnique_invariant_boundary_of_mem_stripParentGroundSpace,
        TNLean.PEPS.finrank_stripParentGroundSpace}
    \leanok
    \uses{thm:peps_g_injective_strip_parent_space}
    Under the preceding hypotheses, each vector in $\mathcal K_n$
    is produced by a unique invariant virtual boundary operator.
    Moreover
    \[
        \dim\mathcal K_n=\dim\operatorname{End}(\mathcal V)^G.
    \]
\end{corollary}

\begin{proof}\leanok
    \uses{thm:peps_g_injective_invariants_range}
    The full strip tensor is $G$-injective by concatenation. Its
    boundary map identifies the invariant virtual operator space
    with its range, which is the common adjacent-pair space.
\end{proof}
